\chapter*{Lời Mở Đầu}
\addcontentsline{toc}{chapter}{Lời Mở Đầu}
\markboth{Lời Mở Đầu}{Lời Mở Đầu}

\begin{truyen}{Lời Mở Đầu}{Opening Words}
\chuhoa{C}{uốn sách này đi qua những nghịch lý rất gần với đời sống: được mà cũng là mất, thắng mà chưa chắc đã hơn, đúng mà đôi khi vẫn thiếu lòng người.}
\textit{(This book moves through paradoxes very close to life: gains that are also losses, victories that may not truly elevate us, and truths that can still lack human warmth.)}

Cuốn này gồm ba phần lớn: \textbf{Được và Mất}, \textbf{Thắng và Thua}, \textbf{Đúng và Sai}. Từ ba phần ấy, sách mở rộng thành 20 chương và 100 truyện ngắn, chạm tới thêm những cặp nghịch lý khác như \textbf{Có và Không}, \textbf{May và Rủi}, \textbf{Nhanh và Chậm}, \textbf{Mạnh và Yếu}, \textbf{Thực và Ảo}.
\textit{(This book contains three major parts: \textbf{Gain and Loss}, \textbf{Winning and Losing}, and \textbf{Right and Wrong}. From those three parts, it expands into 20 chapters and 100 short stories, reaching further paradoxes such as \textbf{Having and Not Having}, \textbf{Luck and Misfortune}, \textbf{Fast and Slow}, \textbf{Strong and Weak}, and \textbf{Real and Illusion}. )}

Nếu đọc xong mà người đọc sống chậm hơn một chút, bớt hơn thua một chút, và thành thật với chính mình thêm một chút, thì cuốn sách này đã làm đúng việc của nó.
\textit{(If after reading, the reader lives a little more slowly, competes a little less, and becomes a little more honest with themselves, then this book has done its work.)}
\end{truyen}

\begin{baihoc}
Một cuốn sách dài hơn chỉ có ý nghĩa khi nó giúp người đọc nhìn đời kỹ hơn và nhìn mình thật hơn.
\textit{(A longer book only matters when it helps readers see life more clearly and see themselves more honestly.)}
\end{baihoc}
\ngancach